% ============================================================
%  01_legislative_brief.tex
%  Autonomy, Capacity, and the Right to Exit
%  LEGISLATIVE BRIEF — For policymakers, staffers, committee members
%  Justin Bogner, 2026
%  Compile: xelatex 01_legislative_brief.tex (twice)
% ============================================================

\documentclass[11pt, letterpaper]{article}
\usepackage{campaign}
\usepackage{xspace}
\usepackage{colortbl}

\campaignheader%
  {Autonomy, Capacity, and the Right to Exit}%
  {Legislative Brief --- Bogner, 2026}

\begin{document}

\campaigntitle%
  {AUTONOMY, CAPACITY, AND THE RIGHT TO EXIT}%
  {A Legislative Brief on Capacity-Based Medical Aid in Dying}%
  {Legislative Brief}%
  {Justin Bogner}%
  {2026}

\suitenote

\begin{policybox}[title={\sffamily\bfseries\small\color{C@accent} Purpose of this Brief}]
This brief presents the case for replacing suffering-based and terminal-diagnosis-based
Medical Aid in Dying (MAiD) eligibility criteria with a \textbf{capacity-based standard}.
It is written for legislative audiences and includes: the policy rationale, a summary of
current jurisdictional models, the proposed eligibility and procedural framework, model
statutory language, and responses to the primary legislative objections.
\end{policybox}

\bigskip

\section{The Policy Problem in Plain Language}

Current MAiD legislation requires one of two things: a terminal diagnosis within a defined
prognosis window, or demonstration of ``unbearable'' suffering. Both criteria share the same
structural flaw: they require the state to evaluate the quality and legitimacy of a person's
subjective experience before recognizing their right to act on it.

This is not the standard medicine applies anywhere else. A patient may refuse dialysis,
chemotherapy, or mechanical ventilation without proving their suffering is ``unbearable.''
The standard is \emph{informed consent}: they understand what they are declining, and they
decline it. That standard should apply consistently to decisions about existence itself.

The practical consequences of the current framework are three:

\begin{enumerate}
  \item Persons with stable, long-considered preferences for death are denied access because
    they cannot demonstrate suffering in a clinically legible form.
  \item The clinical system defaults to coercive retention---hospitalization, safety contracts,
    discharge into unchanged conditions---at significant cost to public health budgets and
    with no demonstrated improvement in outcomes.
  \item Where socioeconomic conditions (inadequate housing, insufficient disability support,
    chronic isolation) are driving MAiD requests, those conditions are neither documented
    nor politically attributed. The state's failure remains invisible.
\end{enumerate}

\section{The Proposed Standard: Decisional Capacity}

The proposed eligibility criterion is \textbf{decisional capacity}. This replaces both the
terminal-diagnosis requirement and the suffering requirement. It does not lower the bar for
access; it changes the nature of the bar from condition-based to comprehension-based---
consistent with informed refusal doctrine already operative in medical law.

\subsection{What Decisional Capacity Requires}

\begin{itemize}
  \item \textbf{Understanding of death.} The applicant understands that death is permanent
    and irreversible.
  \item \textbf{Informed awareness of alternatives.} The applicant has received documented
    presentation of all available alternatives and understands their realistic outcomes.
    \emph{Prior utilization is not required.} Informed refusal, following full disclosure,
    is sufficient and binding.
  \item \textbf{Current cognitive integrity.} The applicant is free from acute psychotic or
    crisis states that materially distort comprehension \emph{at the time of evaluation}.
    A history of psychiatric illness does not disqualify; current incapacity does.
  \item \textbf{Preference stability.} The preference is stable across the mandatory
    assessment period, not reducible to a situational response to a transient condition.
\end{itemize}

\subsection{What It Does Not Require}

\begin{itemize}
  \item A terminal prognosis or defined survival window.
  \item Demonstrated or externally validated suffering.
  \item Prior or current utilization of psychiatric, pharmacological, or other treatment.
  \item Family consent or treating physician agreement.
\end{itemize}

\section{Current Jurisdictional Landscape}

\begin{table}[h!]
\centering
\small
\renewcommand{\arraystretch}{1.45}
\begin{tabularx}{\textwidth}{>{\bfseries\sffamily\small}p{2.6cm} X >{\centering\arraybackslash}p{1.6cm} X}
\toprule
\rowcolor{C@accent}
\thead{Jurisdiction} &
\thead{Eligibility} &
\thead{Terminal Req.} &
\thead{Notes} \\
\midrule
Belgium / Netherlands
  & Unbearable suffering incl.\ psychiatric & No
  & Operative 20+ years; no exponential expansion observed \\
Canada (2026)
  & Grievous \& irremediable; Track~2 non-foreseeable & No (Track~2)
  & 732 Track~2 provisions (2024); 61.5\% disability; 17\% YoY growth \\
Switzerland
  & Any (assisted suicide organizations) & No
  & Existential cases permitted without clinical diagnosis \\
United States (13 states + DC)
  & Terminal $\leq$6 months & Yes
  & Most restrictive model; no non-terminal eligibility \\
\textbf{Proposed}
  & \textbf{Decisional capacity only} & \textbf{No}
  & \textbf{Consistent with informed refusal doctrine} \\
\bottomrule
\end{tabularx}
\caption{MAiD Eligibility Models, 2026. Source: Health Canada Sixth Annual Report (2025).}
\end{table}

\begin{findingbox}[title={\sffamily\bfseries\small\color{C@accent} Key Data Point}]
Belgium and the Netherlands have operated under broad eligibility criteria---including
psychiatric cases---for over two decades. \textbf{Neither jurisdiction has experienced
exponential expansion or collapse of procedural safeguards.} The longitudinal evidence
does not support the slippery slope prediction.
\end{findingbox}

\section{The Procedural Safeguards}

A capacity-based standard is not an unguarded standard. The architecture below
distinguishes resolved autonomous preference from crisis and protects against coercion.

\subsection{Two-Stage Assessment Protocol}

\textbf{Stage 1 --- decisional capacity} (criteria above). Assessed across multiple
independent evaluations by clinicians unaffiliated with the treating team.

\textbf{Stage 2 --- Process Comprehension.} The applicant must demonstrate understanding of
why the safeguards exist. Diagnostically significant: a person in acute crisis will rarely
tolerate a 90-day procedure with multiple independent evaluations. The process itself
differentiates resolved preference from reactive crisis more reliably than a single interview.

\subsection{Mandatory Waiting Period}

\textbf{90 days minimum} from first application to provision for all non-terminal cases.
Applications may be withdrawn at any point without prejudice, penalty, or adverse effect
on future care access.

\subsection{Independent Coercion Screening}

\begin{itemize}
  \item Confidential interview by an independent social worker or patient advocate
    with no institutional stake in the outcome.
  \item Structured review of financial circumstances and dependency relationships.
  \item Mandatory offer of independent legal assistance for advance planning.
\end{itemize}

\subsection{Mandatory Alternatives Presentation}

Evaluators must present every realistic alternative in sufficient specificity for genuine
comprehension: housing assistance, income support, disability services, palliative care,
psychiatric intervention, peer support. Presentation documented. Referrals offered.
Declinations recorded with applicant's stated reasons.

\subsection{Socioeconomic Documentation Requirement}

Where an applicant identifies inadequate housing, income, or disability supports as a
condition of their preference, the administering authority \textbf{must} formally report
this to the relevant government ministry. This is not optional. It is the mechanism by
which systemic failures become politically legible and attributable.

\subsection{Absolute Revocation}

Unconditional. Available until the final moment. No locked-in provision date permitted.
No explanation required. No adverse consequence.

\section{Principal Legislative Objections}

\begin{objectionbox}[title={\sffamily\bfseries\small\color{C@warning}
  Objection 1: ``This will be abused by people in temporary crisis''}]
\textbf{Response:} The procedural architecture addresses this structurally. A 90-day
mandatory assessment period with multiple independent evaluations is incompatible with
impulsive crisis decision-making. Persons in acute crisis want exit \emph{now}; the
timeline filters that case class. The empirical record from Belgium and the Netherlands
confirms: crisis cases do not dominate uptake under broad eligibility frameworks.
\end{objectionbox}

\medskip

\begin{objectionbox}[title={\sffamily\bfseries\small\color{C@warning}
  Objection 2: ``Vulnerable people will be pressured into it''}]
\textbf{Response:} Coercion is a real concern addressed procedurally, not by
categorical exclusion. Independent screening, financial review, and absolute revocation
create multiple protective layers. Categorical exclusion of vulnerable populations is
itself coercive: it overrides expressed preferences on the supposition that they might
have been influenced---a standard that, applied consistently, would invalidate a
substantial proportion of all medical decision-making.
\end{objectionbox}

\medskip

\begin{objectionbox}[title={\sffamily\bfseries\small\color{C@warning}
  Objection 3: ``The disability community opposes this''}]
\textbf{Response:} The disability rights concern---that expanded MAiD enables the state
to substitute death for adequate support---is taken seriously and addressed structurally.
The mandatory socioeconomic documentation requirement makes state failure publicly
attributable. Where inadequate housing or disability support drives a provision, the
state documents and reports it. This creates political accountability that the current
framework, which processes these deaths quietly, does not.
\end{objectionbox}

\medskip

\begin{objectionbox}[title={\sffamily\bfseries\small\color{C@warning}
  Objection 4 (US): ``\textit{Glucksberg} forecloses the constitutional argument''}]
\textbf{Response:} \textit{Washington v.\ Glucksberg}, 521 U.S.\ 702 (1997), and
\textit{Vacco v.\ Quill}, 521 U.S.\ 793 (1997), are the controlling US constitutional
precedents on assisted dying. They held, respectively, that there is no fundamental
right to physician-assisted suicide under substantive due process, and that the
act/omission distinction survives rational basis review under equal protection.

Neither forecloses this framework. This is primarily a \textit{legislative} proposal,
not a constitutional mandate. No constitutional decision is required for state legislatures
to adopt the capacity standard through Death with Dignity reform. Where a constitutional
argument is advanced, it operates through the state-created-danger doctrine
(\textit{DeShaney v.\ Winnebago County}, 489 U.S.\ 189 (1989)), not as a new substantive
due process claim. The protected interest is bodily integrity---already recognized in
\textit{Cruzan} (1990)---not a newly minted right. And \textit{Vacco}'s rational basis
for the act/omission distinction is undermined by the Causation Fiction: when a stable
patient orders a ventilator removed, the immediate cause of death is the active step of
removal. The factual premise underlying this rational basis is increasingly challenged by clinical reality.
\end{objectionbox}

\medskip

\begin{objectionbox}[title={\sffamily\bfseries\small\color{C@warning}
  Objection 5: ``It contradicts the sanctity of life''}]
\textbf{Response:} Medical law already recognizes that competent adults may refuse
life-sustaining treatment---ventilators, dialysis, nutrition---without this being
construed as a violation of the sanctity of life. This framework extends the same
principle consistently. The distinction the current law draws between death by treatment
refusal and death by self-directed action is historically contingent, not philosophically
necessary.
\end{objectionbox}

\section{Model Statutory Language}

\begin{statutebox}
\sffamily\small\setlength{\parskip}{0.55em}

\textbf{\S\,1 --- Eligible Person.}
Any adult who demonstrates decisional capacity as assessed under this Act.
Terminal diagnosis and proof of suffering are not required.

\textbf{\S\,2 --- decisional capacity.}
The demonstrated ability to: (a)~understand the permanent and irreversible nature of
death; (b)~demonstrate informed awareness of available alternatives including medical,
psychiatric, financial, and social supports and their realistic outcomes; (c)~express a
preference stable across the assessment period; (d)~be free from acute psychotic or
crisis states that materially distort comprehension at the time of evaluation.
A history of psychiatric diagnosis is not a disqualifier.

\textbf{\S\,3 --- Informed Refusal.}
An applicant is not required to have utilized any alternative prior to application.
Documented informed awareness of alternatives and reasons for declining them shall
satisfy this element. The evaluating clinician shall assess comprehension, not
compliance history.

\textbf{\S\,4 --- Assessment Period.}
Minimum 90 days from first application for all non-terminal provisions.
Multiple independent evaluations required. Applications may be withdrawn at any time
without prejudice, penalty, or adverse effect on future access to services.

\textbf{\S\,5 --- Coercion Screening.}
Every application shall include: (a)~confidential assessment by an independent social
worker or patient advocate unaffiliated with the treating team; (b)~structured review of
financial and dependency circumstances; (c)~mandatory offer of independent legal assistance.

\textbf{\S\,6 --- Socioeconomic Documentation.}
Where an applicant identifies inadequate housing, income, disability support, or social
connection as a primary or contributing basis for their preference, the administering
authority shall transmit a formal written report to the relevant government ministry
documenting these conditions as a driver of the provision within 30 days.

\textbf{\S\,7 --- Revocation.}
Absolute. Available at any point prior to provision. No explanation required. No adverse
consequence of any kind. No provision date may be locked in advance.

\textbf{\S\,8 --- Conscientious Objection.}
No clinician is required to participate. Objectors must provide timely referral to a
willing provider. Refusal to refer constitutes patient abandonment.

\textbf{\S\,9 --- Annual Reporting.}
The administering authority shall publish an annual public report documenting: total
provisions; demographic breakdown; alternatives presented and accepted or declined;
socioeconomic conditions documented as drivers under \S\,6.

\textbf{\S\,10 --- Standard of Proof.}
Decisional capacity must be established by clear and convincing evidence across the full
assessment period. This standard applies to each of the four functional elements of the
decisional capacity definition. The burden of establishing capacity rests on the assessment
record; the applicant is not required to disprove incapacity.

\textbf{\S\,11 --- Evaluator Disagreement.}
Where independent evaluators reach divergent conclusions regarding the presence of
decisional capacity, the matter shall be referred to an independent Clinical Ethics Committee
within fourteen days. A finding of capacity requires reasoned concurrence by at least two
independent evaluators. Where the Clinical Ethics Committee is unable to resolve the
disagreement, the matter shall be referred to the administrative review panel established
under \S\,12.

\textbf{\S\,12 --- Review and Appeal.}
Any applicant whose request is denied following the Stage~1 assessment shall have the
right to: (a)~request a fresh independent evaluation by a new panel within 30 days;
(b)~seek expedited review by an administrative tribunal within 60 days of any adverse
panel decision; and (c)~pursue judicial review of any adverse tribunal decision. All
denial decisions must be issued in writing with reasoned findings. The administering
authority bears the burden of justifying any denial on review.
\end{statutebox}

\section{Fiscal and Political Framing}

\subsection{Cost Considerations}

Coercive psychiatric retention carries substantial public cost: involuntary hospitalization,
emergency department cycling, long-term institutional placement. A capacity-based MAiD
framework reduces administrative burden on crisis infrastructure for the population it
recognizes, while the mandatory documentation requirement creates financial pressure on
governments to invest in the support alternatives that reduce MAiD uptake. The fiscal
incentives align with the ethical ones.

\subsection{Political Framing}

Polling consistently shows majority public support for expanded end-of-life choice when
framed around autonomy rather than specific conditions. The international comparators reduce
the ``extreme'' characterization: Belgium, the Netherlands, Canada, and Switzerland have all
implemented versions of expanded eligibility without predicted social harm. The policy space
is occupied by functioning democracies. Characterizing capacity-based eligibility as
unprecedented is no longer accurate.

\section*{Summary of Legislative Recommendations}

\begin{enumerate}
  \item Replace terminal-diagnosis and suffering-based eligibility with a
    \textbf{decisional capacity standard}, consistent with informed refusal doctrine.
  \item Implement a \textbf{two-stage assessment protocol} with a 90-day minimum
    waiting period for non-terminal cases.
  \item Require \textbf{mandatory documented alternatives presentation}.
  \item Establish \textbf{independent coercion screening} in every assessment.
  \item Mandate \textbf{socioeconomic documentation and ministerial reporting}
    where inadequate support is identified as a driver.
  \item Protect \textbf{absolute revocation rights} until the moment of provision.
  \item Require \textbf{annual public reporting} on provisions and drivers.
\end{enumerate}

\bigskip

{\small\sffamily\color{C@midgray}
\textbf{Key Sources:} Health Canada, \textit{Sixth Annual Report on MAiD in Canada} (2024 data; released November 2025);
\textit{Carter v.\ Canada (AG)}, 2015 SCC 5; Beauchamp \& Childress,
\textit{Principles of Biomedical Ethics} (8th ed.); Government of Canada,
\textit{Bill C-7} (2024).}

\end{document}
